\section*{Sets and Indices}

\begin{itemize}
    \item $Q = \{0,1,2,3\}$: quarters in chronological order, where
    \[
    0=\text{Winter},\quad 1=\text{Spring},\quad 2=\text{Summer},\quad 3=\text{Autumn}.
    \]
    \item $(i,j) \in A := \{(i,j)\in Q\times Q : i \le j\}$: feasible purchase--sale quarter pairs. Timber bought in quarter $i$ is sold in quarter $j$.
    \item $t \in \{0,1,2\}$: end-of-quarter inventory checkpoints for storage-capacity enforcement.
\end{itemize}

\section*{Parameters}

All volumes are modeled in units of $10^4\text{ m}^3$ and all monetary quantities in units of $10^4$ yuan, matching the code.

\begin{itemize}
    \item $p_i$: purchase price in quarter $i$ (code data: \texttt{purchase\_price[quarters[i]]}).
    \item $r_j$: sale price in quarter $j$ (code data: \texttt{sale\_price[quarters[j]]}).
    \item $d_j$: estimated maximum sales volume in quarter $j$ (code data: \texttt{max\_sales[quarters[j]]}).
    \item $C$: maximum storage capacity (code: \texttt{max\_storage\_wm3}).
    \item $a$: base storage cost per stored unit (code: \texttt{storage\_cost\_a\_w}).
    \item $b$: additional storage cost per stored unit per quarter of storage (code: \texttt{storage\_cost\_b\_w}).
    \item For each $(i,j)\in A$, let
    \[
    u_{ij} := j-i
    \]
    denote the number of quarters timber is stored before sale.
    \item For each $(i,j)\in A$, define the per-unit storage cost used in the code:
    \[
    s_{ij} :=
    \begin{cases}
    0, & j=i,\\
    a + b(j-i), & j>i.
    \end{cases}
    \]
\end{itemize}

\section*{Decision Variables}

\begin{itemize}
    \item $x_{ij} \ge 0$ for $(i,j)\in A$: amount of timber purchased in quarter $i$ and sold in quarter $j$ (code: \texttt{x[i,j]}).
\end{itemize}

\section*{Objective}

Maximize total profit:
\[
\max \sum_{(i,j)\in A} \bigl(r_j - p_i - s_{ij}\bigr)x_{ij}.
\]

Equivalently, using the code's piecewise storage-cost logic,
\[
\max \left[
\sum_{i\in Q} (r_i-p_i)x_{ii}
+
\sum_{\substack{(i,j)\in A\\ j>i}}
\bigl(r_j-p_i-a-b(j-i)\bigr)x_{ij}
\right].
\]

\section*{Constraints}

\begin{itemize}
    \item \textbf{Quarterly sales limits.} For each sale quarter $j\in Q$,
    \[
    \sum_{i=0}^{j} x_{ij} \le d_j.
    \]
    
    \item \textbf{End-of-quarter storage capacity.} For each $t\in\{0,1,2\}$, the amount purchased by or before quarter $t$ and sold after quarter $t$ cannot exceed capacity:
    \[
    \sum_{i=0}^{t}\sum_{j=t+1}^{3} x_{ij} \le C.
    \]
    
    \item \textbf{Nonnegativity.}
    \[
    x_{ij} \ge 0 \qquad \forall (i,j)\in A.
    \]
\end{itemize}

\section*{Notes on Implementation}

\begin{itemize}
    \item The model is implemented as a linear program and solved with \texttt{GUROBI\_CMD} through the PuLP-compatible interface.
    \item There are no purchase-supply limits in the code. Procurement is restricted only indirectly through sales limits and storage-capacity constraints.
    \item The model enforces that all timber must be sold no later than Autumn by defining variables only for pairs with $j\le 3$.
    \item Storage cost is charged only when $j>i$. Same-quarter purchase and sale ($j=i$) incur no storage cost.
    \item The code converts the warehouse capacity from $\text{m}^3$ to $10^4\text{ m}^3$ via \texttt{max\_storage\_wm3 = max\_storage / 10000}.
    \item The code comments indicate a unit conversion of storage costs, but the implemented coefficients used in optimization are numerically
    \[
    a = 70,\qquad b = 100,
    \]
    in units of $10^4$ yuan per $10^4\text{ m}^3$ and per quarter, respectively, because
    \texttt{storage\_cost\_a\_w = storage\_cost\_a * 10000 / 10000} and similarly for \texttt{storage\_cost\_b\_w}. The formulation above matches these implemented coefficients exactly.
\end{itemize}
